\documentclass{article}
\usepackage[preprint]{neurips_2026}
\PassOptionsToPackage{numbers}{natbib}
% \usepackage[margin=1in]{geometry}
% \setlength{\columnsep}{0.25in}
\usepackage[utf8]{inputenc}
\usepackage[T1]{fontenc}
\usepackage{hyperref}
\usepackage{url}
\usepackage{booktabs}
\usepackage{multirow}
\usepackage{array}
\usepackage{tabularx}
\usepackage{enumitem}
\usepackage{amsmath,amssymb,amsthm}
\usepackage{amsfonts}
\usepackage{nicefrac}
\usepackage{microtype}
\usepackage{xcolor}
\usepackage{graphicx}
\usepackage{algorithm}
\usepackage{algpseudocode}
\usepackage{caption}
\usepackage{makecell}

\newtheorem{theorem}{Theorem}
\newtheorem{proposition}{Proposition}
\newtheorem{lemma}{Lemma}
\newtheorem{assumption}{Assumption}
\newtheorem{corollary}{Corollary}

\newcommand{\method}{\textsc{EvoRank-LoRA}}
\newcommand{\Rmax}{R_{\max}}
\newcommand{\TODO}{\textit{TBD}}

\title{\method: Evolutionary Bidirectional Rank Adaptation for Low-Rank Fine-Tuning}

\author{%
  Anonymous Authors \\
  Anonymous Institution \\
  \texttt{anonymous@anonymous.edu}
}

\begin{document}
\maketitle

\begin{abstract}
LoRA is an effective parameter-efficient fine-tuning method, but it typically uses a fixed rank throughout training. This is restrictive: rank demand varies across layers and over time. We propose \method, which treats rank allocation as a discrete structural optimization problem and solves it with an evolution strategy (ES). Each LoRA layer is parameterized in a maximum-rank super-space, and a binary mask selects the active rank-1 components. Training alternates between gradient updates of the active LoRA weights and ES-based structure updates that expand, reduce, or reallocate rank across layers. We show that expansion enlarges the feasible update space, reduction incurs a controlled local loss increase under smoothness assumptions, elitist ES improves the regularized structural objective during structure selection, and the full alternating procedure converges to a block-wise local optimum under a standard two-timescale view. \method is simple, modular, and compatible with standard transformer fine-tuning pipelines.
\end{abstract}

\section{Introduction}
Parameter-efficient fine-tuning (PEFT) adapts large pre-trained models by training only a small set of additional parameters. Among PEFT methods, LoRA is particularly effective: it freezes the backbone and learns low-rank updates for selected linear projections \citep{hu2022lora}. In practice, however, LoRA still relies on a manually chosen rank, and that rank is usually shared across all injected layers.

A fixed rank is convenient but restrictive. Transformer layers do not contribute equally to adaptation, and their capacity requirements need not remain constant during training. A small rank may underfit important layers, while a large rank may waste parameters on layers that need little adaptation. This motivates a dynamic alternative in which rank is allocated where it is useful instead of being fixed in advance.

The difficulty is that rank adaptation is discrete. Gradient descent can optimize the continuous LoRA weights, but it does not naturally decide when to activate a new rank-1 component, remove a redundant one, or move budget across layers. We address this with \method, which combines gradient-based LoRA training with an outer-loop evolution strategy (ES) over rank structure.

The method places each LoRA layer in a maximum-rank super-space and uses a binary mask to select the active rank-1 components. The inner loop updates the active LoRA parameters. The outer loop mutates the structure through expansion, reduction, and cross-layer reallocation, and retains only improving candidates. This yields a simple bidirectional mechanism that can grow capacity when needed and remove it when redundant.

Our contributions are threefold. First, we propose an ES-based LoRA framework with dynamic bidirectional rank adaptation. Second, we introduce a rank super-space parameterization with output-stabilized rsLoRA scaling and structure-aware optimizer-state resets for online rank mutation. Third, we provide theoretical support for the core operators and a submission-ready experimental protocol for evaluating the method.


\section{Related Work}
\label{sec:related}

\paragraph{PEFT.}
PEFT methods adapt frozen pre-trained models with a small number of trainable parameters, including adapters \citep{houlsby2019adapter}, prefix tuning \citep{li2021prefix}, prompt tuning \citep{lester2021power}, and IA$^3$ \citep{liu2022fewshot}. LoRA is a particularly simple and strong instance of this idea.

\paragraph{LoRA variants.}
LoRA represents weight updates as low-rank factors \citep{hu2022lora}. Follow-up work studies adaptive budget allocation \citep{zhang2023adalora}, weight decomposition \citep{liu2024dora}, and other structured variants. Our work is closest to adaptive-rank LoRA, but focuses on explicit discrete rank evolution during training.

\paragraph{Structure adaptation and ES.}
Dynamic structure adaptation is related to pruning, low-rank compression, and neural architecture search. Evolution strategies are attractive in this setting because they can optimize discrete design variables without gradients \citep{rechenberg1973evolutionsstrategie,schwefel1977numerische,salimans2017es,real2019regularized}. We use ES only for outer-loop rank-structure search, while continuous LoRA parameters are still trained by gradient descent.

\section{Method}
\label{sec:method}

\subsection{Preliminaries: LoRA}
For a pre-trained weight matrix $W \in \mathbb{R}^{d \times k}$, standard LoRA learns a low-rank update
\begin{equation}
    W' = W + \Delta W,
\end{equation}
where
\begin{equation}
    \Delta W = \frac{\alpha}{\sqrt{r}}BA,
\end{equation}
with $B \in \mathbb{R}^{d \times r}$ and $A \in \mathbb{R}^{r \times k}$. Here $r$ denotes the LoRA rank. Our implementation follows an rsLoRA-style parameterization in which the effective update is scaled by $\alpha/\sqrt{r}$ to stabilize output magnitude as the active rank changes. Conventional LoRA uses a fixed rank chosen before training and often shares that rank across all LoRA-injected layers.

\subsection{Overview}
\method replaces the fixed-rank design with alternating continuous and discrete updates:
\begin{enumerate}[leftmargin=1.2em]
    \item \textbf{Parameter step:} update the active LoRA factors by gradient descent.
    \item \textbf{Structure step:} mutate and select the rank structure with ES.
\end{enumerate}
Each LoRA layer is embedded in a maximum-rank super-space, and a binary mask selects the active rank-1 components.

\subsection{Rank Super-Space Parameterization}
For the $\ell$-th LoRA layer, we write
\begin{equation}
    \Delta W_\ell = \frac{\alpha}{\sqrt{r_\ell}}\sum_{i=1}^{\Rmax} m_{\ell,i} b_{\ell,i} a_{\ell,i}^{\top},
\end{equation}
where $m_{\ell,i} \in \{0,1\}$ indicates whether the $i$-th rank-1 component is active. In matrix form,
\begin{equation}
    \Delta W_\ell = \frac{\alpha}{\sqrt{r_\ell}} B_\ell M_\ell A_\ell,
\end{equation}
with
\begin{equation}
\begin{aligned}
    B_\ell &\in \mathbb{R}^{d_\ell \times \Rmax}, \\
    A_\ell &\in \mathbb{R}^{\Rmax \times k_\ell}, \\
    M_\ell &= \mathrm{diag}(m_{\ell,1}, \dots, m_{\ell,\Rmax}).
\end{aligned}
\end{equation}
The effective rank is
\begin{equation}
    r_\ell = \sum_{i=1}^{\Rmax} m_{\ell,i}.
\end{equation}
This parameterization supports both expansion and reduction without rebuilding the module. To preserve the magnitude of the retained active update immediately after a structural mutation, the implementation applies a multiplicative compensation factor $\sqrt{r_{\mathrm{new}}/r_{\mathrm{old}}}$ to the surviving active components after each expansion or reduction step.

\subsection{Optimization Objective}
Let $z$ denote the collection of all structural masks and let $\Theta$ denote all continuous LoRA parameters. The validation-side structural objective is
\begin{equation}
    \mathcal{J}(z) = \mathcal{L}_{\mathrm{val}}(\Theta; z) + \lambda_c C(z),
\end{equation}
where $C(z)$ is a complexity term such as total active rank,
\begin{equation}
    C(z) = \sum_{\ell=1}^{L} r_\ell,
\end{equation}
or a size-aware cost,
\begin{equation}
    C(z) = \sum_{\ell=1}^{L}(d_\ell + k_\ell)r_\ell.
\end{equation}
The inner loop optimizes $\Theta$ on training data, while the outer loop updates $z$ using the regularized validation reward.

\subsection{Expansion and Reduction Criteria}
We quantify whether a layer needs more rank using the idealized demand signal
\begin{equation}
    g_\ell = \left\|\frac{\partial \mathcal{L}}{\partial \Delta W_\ell}\right\|_F.
\end{equation}
In practice, materializing the full $d_\ell \times k_\ell$ gradient matrix is prohibitively expensive. The implementation therefore uses the surrogate
\begin{equation}
    \underline{g}_\ell
    =
    \left\|
    \Pi_{\mathrm{row}(A_\ell)}
    \left(
    \frac{\partial \mathcal{L}}{\partial \Delta W_\ell}
    \right)
    \right\|_F
    \le g_\ell,
\end{equation}
which reconstructs the gradient only in the row space of the active $A_\ell$ factors via a pseudo-inverse and requires only small $r_\ell \times r_\ell$ linear solves rather than explicit dense gradient storage. Unless otherwise stated, the implementation-oriented sections use $\underline{g}_\ell$ for controller statistics and for gradient-guided expansion initialization. A large $\underline{g}_\ell$ indicates that the current active low-rank space is likely insufficient. We also define a component importance score
\begin{equation}
    s_{\ell,i} = \|b_{\ell,i}\|_2\,\|a_{\ell,i}\|_2,
\end{equation}
and combine these into an expansion utility
\begin{equation}
    u_\ell = \alpha \underline{g}_\ell + \beta \cdot \frac{1}{\max(r_\ell,1)}\sum_{i:m_{\ell,i}=1} s_{\ell,i}.
\end{equation}
Large $u_\ell$ favors expansion. For reduction, we use
\begin{equation}
    s^{\mathrm{red}}_{\ell,i} = \alpha_1\|b_{\ell,i}\|_2\|a_{\ell,i}\|_2 + \alpha_2\left|\left\langle \frac{\partial \mathcal{L}}{\partial \Delta W_\ell}, b_{\ell,i}a_{\ell,i}^{\top}\right\rangle\right|.
\end{equation}
Low-scoring active components are candidates for removal.

\subsection{ES-Based Structure Evolution}
Given the current structure $z_t$, the ES samples a local neighborhood $\mathcal{N}(z_t)$ with three mutation types:
\begin{itemize}[leftmargin=1.2em]
    \item \textbf{Expansion:} activate a new rank-1 component in a high-demand layer.
    \item \textbf{Reduction:} deactivate a low-importance active component.
    \item \textbf{Reallocation:} reduce one layer and expand another, approximately preserving the global budget.
\end{itemize}
To avoid combinatorial explosion in cross-layer reallocation, the implementation uses a top-k cross strategy: expansion candidates are ordered by layer demand, pruning candidates are ordered by component importance, and only the first $K_{\mathrm{realloc}}$ cross-layer reallocation pairs are evaluated. Unless otherwise stated, we use $K_{\mathrm{realloc}} = 8$.
Each candidate structure $z'$ receives reward
\begin{equation}
    R(z') = -\mathcal{L}_{\mathrm{val}}(\Theta; z') - \lambda_c C(z').
\end{equation}
We use elitist selection,
\begin{equation}
    z_{t+1} = \arg\max_{z \in \{z_t\} \cup \mathcal{N}(z_t)} R(z),
\end{equation}
so the structure changes only when a better candidate is found.

\begin{algorithm}[t]
\caption{Training procedure of \method}
\label{alg:main}
\begin{algorithmic}[1]
\Require Pre-trained model $W$, training set $\mathcal{D}_{\mathrm{tr}}$, validation subset $\mathcal{D}_{\mathrm{val}}$, maximum rank $\Rmax$, initial rank $r_0$, mutation population size $\lambda$, inner steps $K$
\State Initialize LoRA super-space $\{A_\ell, B_\ell\}_{\ell=1}^L$ with maximum rank $\Rmax$
\State Initialize structure $z_0$ by activating the first $r_0$ components in each LoRA layer
\For{$t = 0,1,\dots,T-1$}
    \For{$k = 1,2,\dots,K$}
        \State Sample mini-batch from $\mathcal{D}_{\mathrm{tr}}$
        \State Update active LoRA parameters $\Theta$ by gradient descent under structure $z_t$
    \EndFor
    \State Compute layer demand surrogate $\underline{g}_\ell$ and component scores $s^{\mathrm{red}}_{\ell,i}$
    \For{$j = 1,2,\dots,\lambda$}
        \State Sample candidate $z_t^{(j)}$ via expansion, reduction, or reallocation
        \State Evaluate reward $R\big(z_t^{(j)}\big)$ on $\mathcal{D}_{\mathrm{val}}$
    \EndFor
    \State Set $z_{t+1} = \arg\max_{z \in \{z_t, z_t^{(1)}, \dots, z_t^{(\lambda)}\}} R(z)$
\EndFor
\State \Return trained LoRA parameters $\Theta$ and learned structure $z_T$
\end{algorithmic}
\end{algorithm}

\begin{figure*}[t]
    \centering
    \fbox{%
    \begin{minipage}{0.96\textwidth}
    \vspace{0.35em}
    \textbf{Method figure layout suggestion.} The final camera-ready figure should contain four blocks arranged from left to right: \textbf{(1) LoRA rank super-space}, showing that each injected projection reserves $\Rmax$ rank-1 slots but activates only a subset; \textbf{(2) Inner-loop parameter update}, showing gradient descent over active LoRA components on training mini-batches; \textbf{(3) Statistics extraction}, showing the layer demand surrogate $\underline{g}_\ell$ and reduction score $s^{\mathrm{red}}_{\ell,i}$; and \textbf{(4) ES outer loop}, showing expansion, reduction, and reallocation mutations evaluated on a validation subset with elitist selection. A compact annotation beneath the figure can highlight that the method jointly optimizes continuous LoRA weights and discrete rank structure.
    \vspace{0.35em}
    \end{minipage}}
    \caption{Schematic of \method. The conceptual figure should emphasize the separation between gradient-based parameter learning and ES-based structure evolution. This placeholder box is intentionally text-only so the manuscript remains self-contained; it can later be replaced with a polished vector diagram without changing the caption or surrounding discussion.}
    \label{fig:method}
\end{figure*}

\subsection{Practical Remarks}
New components can be initialized near zero or from a gradient-informed rank-1 approximation. The latter follows Proposition~\ref{prop:best_rank1_direction} but avoids a costly full SVD by using power iteration to estimate the leading singular vectors of the projected gradient operator. When a previously inactive or pruned component is reactivated, its weights are re-initialized rather than inheriting historical values. The implementation also resets the optimizer state associated with that slot before subsequent updates, which avoids stale momentum from destabilizing the current training dynamics. In practice, the outer loop can be run only every few inner updates.


\subsection{Dynamic Thresholding and Implementation Details}
\label{subsec:impl_details}

We do not use fixed absolute thresholds for rank expansion and rank reduction, since the magnitudes of layer-wise gradients and component importances vary substantially across layers, tasks, and training stages. Instead, \method uses relative and dynamically updated thresholds based on normalized statistics. Specifically, we use max-scaling normalization,
\begin{equation}
    \tilde{x} = \frac{x}{\max(x)+\epsilon},
\end{equation}
which keeps all normalized scores non-negative and preserves the interpretability of the weighted combination.

For each LoRA layer $\ell$, we define a capacity score
\begin{equation}
    u_\ell = \alpha \tilde{g}_\ell + \beta \widetilde{\bar{s}}_\ell,
\end{equation}
where $\tilde{g}_\ell$ is the normalized gradient-based demand score and $\widetilde{\bar{s}}_\ell$ is the normalized average importance of active rank-1 components. A large $u_\ell$ indicates that the current rank allocation of layer $\ell$ is likely insufficient.

For each active rank-1 component $(\ell,i)$, we define an importance score
\begin{equation}
    s^{\mathrm{red}}_{\ell,i} = \alpha_1\|b_{\ell,i}\|_2\|a_{\ell,i}\|_2 + \alpha_2\left|\left\langle \frac{\partial \mathcal{L}}{\partial \Delta W_\ell}, b_{\ell,i}a_{\ell,i}^{\top}\right\rangle\right|,
\end{equation}
where small values of $s_{\ell,i}$ indicate potentially redundant components.

To improve robustness, both scores are smoothed by an exponential moving average (EMA):
\begin{equation}
    \hat{u}_\ell^{(t)}
    =
    \rho \hat{u}_\ell^{(t-1)} + (1-\rho) u_\ell^{(t)},
\end{equation}
\begin{equation}
    \hat{s}_{\ell,i}^{(t)}
    =
    \rho \hat{s}_{\ell,i}^{(t-1)} + (1-\rho) s_{\ell,i}^{(t)},
\end{equation}
where $\rho \in [0.8,0.95]$ is the smoothing factor.

During each structural update, the expansion threshold is defined by the upper percentile of the layer-wise EMA scores:
\begin{equation}
    \tau_{\mathrm{grow}}^{(t)}
    =
    \mathrm{Percentile}\bigl(\{\hat{u}_\ell^{(t)}\}_{\ell=1}^{L},\, p_g\bigr),
\end{equation}
where $p_g$ is typically set to $80\%$. A layer is considered eligible for rank expansion if
\begin{equation}
    \hat{u}_\ell^{(t)} > \tau_{\mathrm{grow}}^{(t)}
    \quad\text{and}\quad
    r_\ell < r_{\max}.
\end{equation}

Similarly, the reduction threshold is defined by the lower percentile of active component scores:
\begin{equation}
    \tau_{\mathrm{prune}}^{(t)}
    =
    \mathrm{Percentile}\bigl(\{\hat{s}_{\ell,i}^{(t)}: m_{\ell,i}=1 \text{ and } (\ell,i)\notin \mathcal{C}^{(t)}\},\, p_p\bigr),
\end{equation}
where $p_p$ is typically set to $10\%$. An active component is considered eligible for removal if
\begin{equation}
    \hat{s}_{\ell,i}^{(t)} < \tau_{\mathrm{prune}}^{(t)}
    \quad\text{and}\quad
    r_\ell > r_{\min}.
\end{equation}

To prevent oscillatory grow-prune behavior, we use three stabilizing mechanisms. First, structural evolution is disabled during an initial warmup stage. This warmup addresses a zero-gradient trap caused by standard LoRA initialization: when $B_\ell$ is initialized at or near zero, the gradients of the corresponding $A_\ell$ rows are initially suppressed, which can make early structural statistics unreliable. Second, a candidate expansion or reduction must satisfy the corresponding threshold for multiple consecutive structural updates before it is executed. Third, a cooldown period is imposed after expansion, during which newly activated components cannot be immediately removed; unlike warmup, this mechanism protects freshly introduced components during later training.

Unless otherwise stated, we apply structural evolution every 200 optimization steps, use an EMA factor of 0.9, set the growth and pruning percentiles to $80\%$ and $10\%$, respectively, and use a warmup stage covering the first $10\%$ of total training steps. We set the minimum and maximum per-layer ranks to $r_{\min}=2$ and $r_{\max}=16$, use persistence lengths $H_g=2$ and $H_p=3$, and apply a cooldown period of two structural updates after each expansion.

\begin{algorithm}[t]
\caption{Dynamic thresholding and candidate identification in \method}
\label{alg:dynamic_threshold}
\footnotesize
\begin{algorithmic}[1]
\Require LoRA parameters $\Theta$, structure masks $z$, update interval $T_{\mathrm{es}}$, smoothing factor $\rho$, growth percentile $p_g$, pruning percentile $p_p$, minimum rank $r_{\min}$, maximum rank $r_{\max}$, warmup steps $T_{\mathrm{warm}}$
\State Initialize EMA statistics from the first observed values of $u_\ell$ and $s_{\ell,i}$
\State Initialize persistence counters $c^{\mathrm{grow}}_\ell \leftarrow 0$, $c^{\mathrm{prune}}_{\ell,i} \leftarrow 0$
\State Initialize cooldown flags for newly expanded components
\For{$t = 1, 2, \dots, T$}
    \State Perform standard gradient-based optimization of active LoRA parameters
    \If{$t \le T_{\mathrm{warm}}$}
        \State \textbf{continue}
    \EndIf
    \If{$t \bmod T_{\mathrm{es}} = 0$}
        \For{each LoRA layer $\ell$}
            \State Compute capacity score $u_\ell$
            \State Update EMA: $\hat{u}_\ell \leftarrow \rho \hat{u}_\ell + (1-\rho) u_\ell$
        \EndFor
        \For{each active component $(\ell,i)$}
            \State Compute importance score $s_{\ell,i}$
            \State Update EMA: $\hat{s}_{\ell,i} \leftarrow \rho \hat{s}_{\ell,i} + (1-\rho) s_{\ell,i}$
        \EndFor
        \State Compute expansion threshold $\tau_{\mathrm{grow}} = \mathrm{Percentile}(\{\hat{u}_\ell\}_{\ell=1}^{L}, p_g)$
        \State Compute pruning threshold $\tau_{\mathrm{prune}} = \mathrm{Percentile}(\{\hat{s}_{\ell,i}: m_{\ell,i}=1, (\ell,i)\notin\mathcal{C}\}, p_p)$
        \For{each layer $\ell$}
            \If{$\hat{u}_\ell > \tau_{\mathrm{grow}}$ and $r_\ell < r_{\max}$}
                \State $c^{\mathrm{grow}}_\ell \leftarrow c^{\mathrm{grow}}_\ell + 1$
            \Else
                \State $c^{\mathrm{grow}}_\ell \leftarrow 0$
            \EndIf
        \EndFor
        \For{each active component $(\ell,i)$}
            \If{$\hat{s}_{\ell,i} < \tau_{\mathrm{prune}}$ and $r_\ell > r_{\min}$ and component $(\ell,i)$ is not in cooldown}
                \State $c^{\mathrm{prune}}_{\ell,i} \leftarrow c^{\mathrm{prune}}_{\ell,i} + 1$
            \Else
                \State $c^{\mathrm{prune}}_{\ell,i} \leftarrow 0$
            \EndIf
        \EndFor
        \State Collect expansion candidates $\{ \ell : c^{\mathrm{grow}}_\ell \ge H_g \}$
        \State Collect pruning candidates $\{ (\ell,i) : c^{\mathrm{prune}}_{\ell,i} \ge H_p \}$
        \State Build ES trial mutations from expansion, reduction, and top-k cross reallocation candidates
        \State Pass the candidate set to Algorithm~\ref{alg:main} for validation-side elitist selection
        \State Update cooldown status of expanded components
    \EndIf
\EndFor
\State \Return Updated statistics and candidate sets for ES-based structure selection
\end{algorithmic}
\end{algorithm}

\section{Theoretical Analysis}
\label{sec:theory}

We analyze the LoRA update space under standard smoothness assumptions and focus on the three core operations of \method: expansion, reduction, and ES-based structure updates. For notational simplicity, the theory below omits the implementation-side rsLoRA scaling factor and studies the corresponding unscaled low-rank update geometry; the implementation reinstates the scaling and compensation rules described in Section~\ref{sec:method}.

\subsection{Preliminaries}
For a single LoRA-injected layer, let
\begin{equation}
    \Delta W = \sum_{i=1}^{\Rmax} m_i b_i a_i^\top,
\end{equation}
where $m_i \in \{0,1\}$ indicates whether the $i$-th rank-1 component is active. Let
\begin{equation}
    r = \sum_{i=1}^{\Rmax} m_i
\end{equation}
be the effective rank. Let $\mathcal{L}(\Delta W)$ denote the objective as a function of the LoRA update for this layer with other parameters fixed.

\begin{assumption}
\label{ass:smoothness}
The loss $\mathcal{L}(\Delta W)$ is $\beta$-smooth with respect to $\Delta W$, i.e., for any compatible matrices $X$ and $Y$,
\begin{equation}
\mathcal{L}(Y) \le \mathcal{L}(X) + \langle \nabla \mathcal{L}(X), Y-X \rangle + \frac{\beta}{2}\|Y-X\|_F^2.
\end{equation}
\end{assumption}

\subsection{Effect of Rank Expansion}
\begin{proposition}[Monotonicity with respect to rank expansion]
\label{prop:rank_monotonicity}
Let
\begin{equation}
    F_r^\star = \min_{\mathrm{rank}(\Delta W) \le r} F(\Delta W)
\end{equation}
for an arbitrary objective $F$. Then for any $r \ge 0$,
\begin{equation}
    F_{r+1}^\star \le F_r^\star.
\end{equation}
\end{proposition}

\begin{proof}
The feasible set for rank at most $r$ is a subset of the feasible set for rank at most $r+1$. Therefore, minimizing over the larger set cannot yield a worse optimum.
\end{proof}

This shows that rank expansion can only enlarge the feasible update space.

\subsection{Optimal Direction of Rank Expansion}
Suppose we add a new rank-1 component $\eta ba^\top$ with sufficiently small $\eta > 0$ and unit vectors $a$ and $b$. A first-order approximation gives
\begin{equation}
    \mathcal{L}(\Delta W + \eta ba^\top)
    \approx
    \mathcal{L}(\Delta W) + \eta\langle \nabla\mathcal{L}(\Delta W), ba^\top \rangle.
\end{equation}

\begin{proposition}[Best rank-1 expansion direction]
\label{prop:best_rank1_direction}
Let $G = \nabla \mathcal{L}(\Delta W)$. Then the solution to
\begin{equation}
    \max_{\|a\|_2 = \|b\|_2 = 1} \langle -G, ba^\top \rangle
\end{equation}
is given by the leading left and right singular vectors of $G$, and the maximum value equals $\sigma_1(G)$.
\end{proposition}

\begin{proof}
The claim follows from the variational characterization of the spectral norm.
\end{proof}

Thus, the dominant gradient direction is a natural signal for expansion.

\subsection{Error Bound for Rank Reduction}
Consider removing the $j$-th active component so that
\begin{equation}
    \Delta W^{(-j)} = \Delta W - b_j a_j^\top.
\end{equation}

\begin{theorem}[Loss increase bound for removing one component]
\label{thm:pruning_bound}
Under Assumption~\ref{ass:smoothness},
\begin{equation}
\label{eq:pruning_bound}
\begin{aligned}
    \mathcal{L}(\Delta W^{(-j)}) - \mathcal{L}(\Delta W)
    &\le \left|\langle \nabla\mathcal{L}(\Delta W), b_j a_j^\top \rangle\right| \\
    &\quad + \frac{\beta}{2}\|b_j\|_2^2\|a_j\|_2^2.
\end{aligned}
\end{equation}
\end{theorem}

\begin{proof}
Apply Assumption~\ref{ass:smoothness} with $X = \Delta W$ and $Y = \Delta W^{(-j)}$. Since $Y-X = -b_j a_j^\top$, we obtain
\begin{equation}
    \mathcal{L}(Y) - \mathcal{L}(X)
    \le
    -\langle \nabla\mathcal{L}(X), b_j a_j^\top \rangle
    + \frac{\beta}{2}\|b_j a_j^\top\|_F^2.
\end{equation}
Using $\|b_j a_j^\top\|_F = \|b_j\|_2\|a_j\|_2$ and upper-bounding the linear term by its absolute value yields the result.
\end{proof}

This bound motivates removing components with small gradient interaction and small norm.

\subsection{Spectral Interpretation of Reduction}
\begin{proposition}[Best rank-$r$ approximation error]
\label{prop:eckart_young}
Let $\Delta W^\star$ have singular value decomposition
\begin{equation}
    \Delta W^\star = \sum_{i=1}^{q} \sigma_i u_i v_i^\top,
    \qquad \sigma_1 \ge \sigma_2 \ge \cdots \ge \sigma_q \ge 0.
\end{equation}
Then
\begin{equation}
    \min_{\mathrm{rank}(X) \le r} \|\Delta W^\star - X\|_F^2
    = \sum_{i=r+1}^{q} \sigma_i^2.
\end{equation}
\end{proposition}

\begin{proof}
This is the Eckart--Young--Mirsky theorem.
\end{proof}

A rapidly decaying spectrum therefore favors aggressive reduction.

\subsection{Monotonic Improvement of ES-Based Structure Search}
Define the structural reward
\begin{equation}
    R(z) = -\mathcal{L}_{\mathrm{val}}(\Theta; z) - \lambda_c C(z).
\end{equation}
The ES update uses elitist selection,
\begin{equation}
    z_{t+1} = \arg\max_{z \in \{z_t\}\cup\mathcal{N}(z_t)} R(z).
\end{equation}

\begin{theorem}[Monotonic reward improvement under elitist selection]
\label{thm:elitist_monotonicity}
Assume the continuous parameters $\Theta$ are fixed during the structural selection step. Then
\begin{equation}
    R(z_{t+1}) \ge R(z_t).
\end{equation}
Equivalently, the regularized structural objective $\mathcal{J}(z) = -R(z)$ is monotonically non-increasing during each structure update.
\end{theorem}

\begin{proof}
Since $z_t$ is included in the candidate set, the selected maximizer cannot have smaller reward than the current structure.
\end{proof}



\subsection{Convergence Analysis}
We focus on the natural notion of convergence for a mixed discrete--continuous procedure: block-wise local convergence under a two-timescale view.

Let $z \in \mathcal{Z}$ denote the discrete rank structure and let $\Theta$ denote the continuous LoRA parameters. Consider the regularized objective
\begin{equation}
    F(\Theta, z) = \mathcal{L}(\Theta; z) + \lambda_c C(z),
\end{equation}
and define the value function
\begin{equation}
    \Phi(z) = \inf_{\Theta} F(\Theta, z).
\end{equation}
We make three standard assumptions: (i) the structure space $\mathcal{Z}$ is finite, (ii) for any fixed $z$, the inner optimization produces limit points that are first-order stationary points of $F(\Theta,z)$, and (iii) the outer ES step uses elitist selection and accepts only strict improvements in $\Phi(z)$ over the mutation neighborhood.

\begin{theorem}[Finite stabilization of the outer structure]
\label{thm:outer_stabilization}
Under the assumptions above, every accepted outer-loop structure update strictly decreases $\Phi(z)$. Consequently, the sequence of accepted structures stabilizes after finitely many updates at some $z^\star \in \mathcal{Z}$ satisfying
\begin{equation}
    \Phi(z^\star) \le \Phi(z), \qquad \forall z \in \mathcal{N}(z^\star).
\end{equation}
That is, $z^\star$ is locally optimal in the mutation neighborhood.
\end{theorem}

\begin{proof}[Proof sketch]
Elitist strict-improvement selection makes $\Phi(z)$ decrease whenever the structure changes. Because $\mathcal{Z}$ is finite, there cannot be infinitely many strict decreases. Hence the accepted structure sequence must terminate at some $z^\star$. At termination, no neighbor yields a smaller value of $\Phi$, which is exactly neighborhood local optimality.
\end{proof}

The outer ES loop therefore cannot change the rank pattern indefinitely.

\begin{theorem}[Block-wise local convergence]
\label{thm:block_convergence}
Under the same assumptions, every limit point $(\Theta^\star, z^\star)$ of the alternating optimization in \method is block-wise locally optimal in the following sense:
\begin{enumerate}[leftmargin=1.2em]
    \item $z^\star$ is locally optimal over the mutation neighborhood with respect to $\Phi(z)$;
    \item $\Theta^\star$ is a first-order stationary point of $F(\Theta, z^\star)$.
\end{enumerate}
Equivalently, the algorithm converges to a point that is locally optimal with respect to discrete structure moves and stationary with respect to continuous parameter updates.
\end{theorem}

\begin{proof}[Proof sketch]
By Theorem~\ref{thm:outer_stabilization}, the outer loop stabilizes after finitely many accepted structure changes. After that time, optimization proceeds with fixed structure $z^\star$ only. The assumed fixed-structure convergence guarantee then implies that every limit point of the inner sequence is a first-order stationary point of $F(\Theta, z^\star)$. Combining this stationary condition with the local optimality of $z^\star$ gives the claim.
\end{proof}

This theorem does not imply global optimality; it shows that structure updates eventually stop and the remaining optimization reduces to fixed-structure LoRA training.

\begin{corollary}[Fixed-structure stationary-point convergence]
\label{cor:fixed_structure_sgd}
Suppose that for a fixed structure $z$, the objective $F(\Theta, z)$ is lower bounded and has Lipschitz-continuous gradients, and that stochastic gradients are unbiased with bounded variance. If the inner loop uses a Robbins--Monro step-size schedule, then
\begin{equation}
    \lim_{T \to \infty} \frac{1}{T} \sum_{t=1}^{T} \mathbb{E}\bigl[\|\nabla_{\Theta} F(\Theta_t, z)\|_2^2\bigr] = 0.
\end{equation}
Hence, once the outer structure has stabilized, the remaining optimization converges to a first-order stationary point in the standard non-convex stochastic optimization sense.
\end{corollary}

\begin{proof}[Proof sketch]
After outer-loop stabilization, the problem reduces to stochastic optimization over the continuous variables $\Theta$ with fixed architecture. The claim then follows from the standard non-convex SGD convergence result under smoothness, bounded-variance noise, and diminishing step sizes.
\end{proof}

\paragraph{Interpretation.}
The convergence result is intentionally local: the outer ES loop stabilizes in finitely many accepted steps, and the inner optimizer converges to a stationary point under the learned structure.

\section{Experiments}
\label{sec:exp}

This section gives a compact experimental scaffold. All \TODO{} entries should be replaced with measured values.

\subsection{Experimental Questions}
We recommend four questions: (1) does bidirectional rank evolution outperform fixed-rank LoRA under comparable budgets; (2) does \method improve the performance--efficiency trade-off relative to AdaLoRA; (3) are both expansion and reduction necessary; and (4) does ES outperform random or greedy structure search.

\subsection{Benchmarks and Backbones}
A strong evaluation should include both discriminative and generative tasks.
\paragraph{Language understanding.}
Use GLUE-style classification tasks such as SST-2, QNLI, and MNLI.
\paragraph{Sequence generation.}
Use summarization or instruction-style generation tasks such as CNN/DailyMail and SAMSum.
\paragraph{Backbones.}
A representative setup is RoBERTa-base for understanding, T5-base for encoder--decoder generation, and one decoder-only model such as LLaMA-2-7B or Mistral-7B. Replace this with the exact models used.

\subsection{Baselines}
Recommended baselines are full fine-tuning, fixed-rank LoRA with $r \in \{4,8,16,32\}$, AdaLoRA, a random-mutation variant of \method, and a greedy variant without ES.

\subsection{Evaluation Metrics}
Report both task quality and efficiency, including accuracy/F1/Rouge-L/BLEU/perplexity as appropriate, along with trainable parameters, average effective rank, peak memory, and training-time overhead.

\subsection{Implementation Details Template}
Table~\ref{tab:hparams} summarizes the implementation details reported in our default configuration and highlights the settings that should still be specified task by task.

\begin{table}[t]
    \centering
    \caption{Recommended implementation details to report. The dynamic-threshold settings below match the default configuration in Section~\ref{subsec:impl_details}. Replace remaining \TODO{} entries with task-specific settings.}
    \label{tab:hparams}
    \begin{tabularx}{\columnwidth}{>{\raggedright\arraybackslash}p{0.40\columnwidth}X}
        \toprule
        Item & Reported setting \\
        \midrule
        LoRA target modules & \TODO{} \\
        Maximum rank $\Rmax$ & 16 \\
        Initial rank $r_0$ & \TODO{} \\
        Minimum rank $r_{\min}$ & 2 \\
        rsLoRA scaling rule & $\alpha / \sqrt{r_\ell}$ with mutation compensation $\sqrt{r_{\mathrm{new}}/r_{\mathrm{old}}}$ \\
        ES population size $\lambda$ & \TODO{} \\
        Structure update interval $T_{\mathrm{es}}$ & 200 optimization steps \\
        Warmup ratio & first 10\% of total steps \\
        EMA factor $\rho$ & 0.9 \\
        Normalization rule & max-scaling $\tilde{x}=x/(\max(x)+\epsilon)$ \\
        Growth percentile $p_g$ & 80\% \\
        Pruning percentile $p_p$ & 10\% \\
        Grow / prune persistence & $H_g=2$, $H_p=3$ \\
        Cooldown after expansion & 2 structural updates \\
        Reallocation strategy & top-k cross with $K_{\mathrm{realloc}}=8$ \\
        Reactivation behavior & re-initialize component weights; reset optimizer state for the affected slot \\
        Reward definition & $-\mathcal{L}_{\mathrm{val}} - \lambda_c C(z)$ \\
        Validation subset size & \TODO{} \\
        Optimizer and LR & \TODO{} \\
        Batch size / epochs & \TODO{} \\
        Expansion initialization & zero-init or gradient-informed via power iteration \\
        \bottomrule
    \end{tabularx}
\end{table}

\subsection{Main Results Table Template}
Table~\ref{tab:main_results} gives a polished main-result template suitable for the final paper.

\begin{table*}[t]
    \centering
    \small
    \caption{Main results template. Replace all \TODO{} entries with measured values. For fair comparison, either match trainable parameter budgets or explicitly report budget differences.}
    \label{tab:main_results}
    \begin{tabularx}{\textwidth}{lcccccc>{\centering\arraybackslash}p{0.12\textwidth}}
        \toprule
        Method & SST-2 & QNLI & MNLI-m & CNN/DM & Avg. & Params & Avg. rank \\
        \midrule
        Full fine-tuning & \TODO{} & \TODO{} & \TODO{} & \TODO{} & \TODO{} & \TODO{} & -- \\
        LoRA-$r=4$ & \TODO{} & \TODO{} & \TODO{} & \TODO{} & \TODO{} & \TODO{} & 4 \\
        LoRA-$r=8$ & \TODO{} & \TODO{} & \TODO{} & \TODO{} & \TODO{} & \TODO{} & 8 \\
        LoRA-$r=16$ & \TODO{} & \TODO{} & \TODO{} & \TODO{} & \TODO{} & \TODO{} & 16 \\
        AdaLoRA & \TODO{} & \TODO{} & \TODO{} & \TODO{} & \TODO{} & \TODO{} & \TODO{} \\
        \method & \TODO{} & \TODO{} & \TODO{} & \TODO{} & \TODO{} & \TODO{} & \TODO{} \\
        \bottomrule
    \end{tabularx}
\end{table*}

\subsection{Ablation Study Template}
Table~\ref{tab:ablation} isolates the contribution of each structural component.

\begin{table}[t]
    \centering
    \small
    \caption{Ablation template for \method. Compare the full method with one-component removals and search variants.}
    \label{tab:ablation}
    \begin{tabularx}{\columnwidth}{>{\raggedright\arraybackslash}Xcc}
        \toprule
        Variant & Metric & $\Delta$ \\
        \midrule
        Full \method & \TODO{} & 0.0 \\
        w/o expansion & \TODO{} & \TODO{} \\
        w/o reduction & \TODO{} & \TODO{} \\
        w/o reallocation & \TODO{} & \TODO{} \\
        Random mutation & \TODO{} & \TODO{} \\
        Greedy mutation & \TODO{} & \TODO{} \\
        \bottomrule
    \end{tabularx}
\end{table}

\subsection{Efficiency Reporting Template}
A separate efficiency table is useful for reporting resource overhead.

\begin{table}[t]
    \centering
    \small
    \caption{Efficiency reporting template. Relative wall-clock time can be normalized to fixed-rank LoRA-$r=8$.}
    \label{tab:efficiency}
    \begin{tabularx}{\columnwidth}{>{\raggedright\arraybackslash}Xcc}
        \toprule
        Method & Rel. time & Params \\
        \midrule
        LoRA-$r=8$ & 1.00$\times$ & \TODO{} \\
        AdaLoRA & \TODO{} & \TODO{} \\
        \method & \TODO{} & \TODO{} \\
        \bottomrule
    \end{tabularx}
\end{table}

\subsection{Qualitative and Diagnostic Analysis}
Useful diagnostics include the final layer-wise rank profile, the evolution of total active rank during training, and a Pareto plot of task metric versus trainable parameters for fixed-rank LoRA, AdaLoRA, and \method.

\section{Limitations}
This draft does not report measured results. A complete submission must evaluate at least one encoder-style and one generation-style backbone, quantify ES overhead, and compare against a strong adaptive-rank baseline. The theory also analyzes the LoRA update space rather than the full non-convex dynamics of large-model fine-tuning.

\section{Conclusion}
We presented \method, which treats LoRA rank allocation as an ES-driven structure search problem. The method combines gradient-based updates of continuous LoRA weights with bidirectional evolution over discrete rank variables. The resulting framework is simple, modular, and provides a practical alternative to fixed-rank LoRA.

\paragraph{Before submission.}
Replace all \TODO{} entries with measured values, convert Figure~\ref{fig:method} into a polished diagram, and update the bibliography with the exact cited versions.

\bibliographystyle{plainnat}
\bibliography{references}
\end{document}
